\section{把癌症重写成世界模型问题}

前两章仍然适用于机器人、视频或语音。本章开始说明为什么肿瘤生物学不是随意挑选的应用：它既是一个多尺度动态系统，又拥有多个互补但昂贵的传感器。临床真正关心的不是“这张图属于哪类”，而是“给这个患者某个治疗后，肿瘤会不会发生有利变化”。

\subsection{宏观世界与肿瘤世界的对应}

从街景到肿瘤，结构保持不变：现实状态不可完全观察，传感器给出部分证据，行动改变未来，查询决定模拟目标。不同之处在于肿瘤的状态变量横跨患者、器官、组织、细胞和分子多个尺度。下面通过宏观世界、免疫微环境和患者模拟器三张图，逐步把抽象组件映射到可测量对象与临床问题。

\lecturefigure{slide-092-macroscopic-world-model.jpg}{宏观世界模型：摄像头、麦克风和文本共同回答行动查询}{00:22:46}

这张回顾页提醒我们，world model 的价值来自闭环：传感器不是为了收集越多越好，而是为了回答行动相关查询。若一个模态昂贵但不会改变治疗排序，它的边际价值可能很低；若一个稀有测量能区分两个治疗响应完全相反的亚群，它就非常关键。

\lecturefigure{slide-094-cancer-immunology.jpg}{一页癌症免疫学：肿瘤与免疫系统在微环境中持续互动}{00:25:04}

图中彩色细胞不是装饰。Tumor microenvironment（肿瘤微环境）包含肿瘤细胞、T 细胞、B 细胞、巨噬细胞和基质等；免疫系统既可能识别并清除异常细胞，也可能被肿瘤抑制、排斥或重编程。治疗响应因此依赖细胞类型、空间位置、状态和信号通路的组合。

\begin{knowledgebox}{第一次出现：肿瘤免疫治疗}
Cancer immunotherapy（癌症免疫治疗）不是直接毒杀所有肿瘤细胞，而是尝试解除免疫抑制、增强识别或改变免疫细胞功能。相同药物在不同患者上反应差异巨大，正是患者条件世界模型想解释的对象。
\end{knowledgebox}

\lecturefigure{slide-095-patient-simulator.jpg}{理想目标：从患者组织出发模拟治疗后的组织级响应}{00:25:20}

图中把患者组织送入 world model，再询问“给药后肿瘤会消退吗”。可形式化为
\[
p_\theta(R\mid X_{\text{patient}},a_{\text{drug}}),
\]
其中 $R$ 是响应，$X_{\text{patient}}$ 是患者多模态状态，$a_{\text{drug}}$ 是药物干预。这个公式看似简单，真正困难的是 $X$ 不可完整测量、药物作用时间很长、训练数据缺少随机干预，而且临床响应标签稀缺。

\lecturefigure{slide-096-dataset-scale.jpg}{讲座时点的数据规模：患者、切片、像素与空间测量共同增长}{00:25:21}

数字页必须按时间理解：它描述 2025 年 5 月讲座时的研究快照，不是 2026 年现状。更重要的不是具体计数，而是层级差异：约千级患者、千级切片、海量像素和空间点位意味着“样本量”不能只按患者数或 token 数单独评价。患者独立性决定泛化，细胞与像素决定局部统计精度。

\teachervoice{00:22:46--00:25:23，老师把问题压缩成一句临床查询：如果把某个治疗给这个患者，肿瘤会不会消退？这比“做一个癌症基础模型”更能约束数据与验证。}

\begin{warningbox}{世界模型不等于临床决策系统}
即使模型能预测组织变化，仍需考虑毒性、剂量、药代动力学、患者合并症和不确定性。组织级模拟是候选证据，不是直接处方。讲义中的“药物发现”应理解为发现与排序假设的一段流程。
\end{warningbox}

\subsection{本章小结}

肿瘤生物学可被重写为一个部分可观测、可干预、多尺度的动态系统。真正目标是患者条件下的治疗响应，而不是静态分类。这个目标要求模型同时处理互补模态、空间结构、有限患者数量和极强的因果验证需求。

